\expandafter\gdef\csname SNside1\endcsname{\node[anchor=north west,inner sep=0pt,text width=48mm,font=\SideFont,align=left] at ([xshift=138mm,yshift=-500.62pt]current page.north west) {\textcolor{SNink}{\textbf{错项排除 · 唤醒与分派是两步}}\par\vspace{1.5mm}I/O完成把阻塞者所缺的事件补齐，通常只使其进入就绪队列。处理器仍可由原进程或另一个高优先级者使用，因此``设备一完成，该进程立即运行''遗漏了调度决定。画状态图时应保留阻塞到就绪、就绪到运行两条边。};}
\expandafter\gdef\csname SNside2\endcsname{\node[anchor=north west,inner sep=0pt,text width=48mm,font=\SideFont,align=left] at ([xshift=138mm,yshift=-509.29pt]current page.north west) {\textcolor{SNink}{\textbf{课程连接 · 模型决定多核并行度}}\par\vspace{1.5mm}多个用户线程若全部复用一个内核执行实体，内核无法同时把它们安排到两个CPU上。若映射到多个内核线程，才可能获得多核并行。判断时把用户线程数、承载实体数、物理核数分开，不能仅凭线程名字推结论。};}
\expandafter\gdef\csname SNside3\endcsname{\node[anchor=north west,inner sep=0pt,text width=48mm,font=\SideFont,align=left] at ([xshift=138mm,yshift=-95.00pt]current page.north west) {\textcolor{SNink}{\textbf{错项排除 · 共享内存只提供共同位置}}\par\vspace{1.5mm}两个进程映射同一物理页后，写入可被对方访问，但这并未规定谁先写、何时读，也未把读改写自动变成原子操作。通信方式解决数据放在哪里，同步机制解决操作怎样协调，二者承担的任务不同。};}
\expandafter\gdef\csname SNside4\endcsname{\node[anchor=north west,inner sep=0pt,text width=48mm,font=\SideFont,align=left] at ([xshift=138mm,yshift=-95.00pt]current page.north west) {\textcolor{SNink}{\textbf{不变量 · 等待时间要扣对部分}}\par\vspace{1.5mm}无I/O的单CPU段模型满足周转＝就绪等待＋服务；若还有阻塞，必须再加阻塞时间。始终把一个进程的整个生存区间按状态分段相加，再套公式，能避免把等待设备的时间全部算进就绪队列等待。};}
\expandafter\gdef\csname SNside5\endcsname{\node[anchor=north west,inner sep=0pt,text width=48mm,font=\SideFont,align=left] at ([xshift=138mm,yshift=-95.00pt]current page.north west) {\textcolor{SNink}{\textbf{最小反例 · 不能挑未来的短任务}}\par\vspace{1.5mm}在时刻2选择作业，时刻3才到达的短作业即使只需1毫秒也不在候选集合。算法先受``已到达且就绪''的约束，随后才比较运行长度或响应比；跳过集合筛选，会把最短优先错误变成预知未来的调度。};}
\expandafter\gdef\csname SNside6\endcsname{\node[anchor=north west,inner sep=0pt,text width=48mm,font=\SideFont,align=left] at ([xshift=138mm,yshift=-180.56pt]current page.north west) {\textcolor{SNink}{\textbf{错项排除 · 共享内存只提供共同位置}}\par\vspace{1.5mm}两个进程映射同一物理页后，写入可被对方访问，但这并未规定谁先写、何时读，也未把读改写自动变成原子操作。通信方式解决数据放在哪里，同步机制解决操作怎样协调，二者承担的任务不同。};}
\expandafter\gdef\csname SNside7\endcsname{\node[anchor=north west,inner sep=0pt,text width=48mm,font=\SideFont,align=left] at ([xshift=138mm,yshift=-95.00pt]current page.north west) {\textcolor{SNink}{\textbf{不变量 · 等待时间要扣对部分}}\par\vspace{1.5mm}无I/O的单CPU段模型满足周转＝就绪等待＋服务；若还有阻塞，必须再加阻塞时间。始终把一个进程的整个生存区间按状态分段相加，再套公式，能避免把等待设备的时间全部算进就绪队列等待。};}
\expandafter\gdef\csname SNside8\endcsname{\node[anchor=north west,inner sep=0pt,text width=48mm,font=\SideFont,align=left] at ([xshift=138mm,yshift=-95.00pt]current page.north west) {\textcolor{SNink}{\textbf{课程连接 · 迁移与Cache局部性相互影响}}\par\vspace{1.5mm}把线程移到空闲核可减少排队，却可能丢掉此前核的Cache与TLB热状态。组成原理的命中率和操作系统的负载均衡在这里共同影响总时间；每核任务数量相等也不保证剩余工作量相等。};}
\expandafter\gdef\csname SNside9\endcsname{\node[anchor=north west,inner sep=0pt,text width=48mm,font=\SideFont,align=left] at ([xshift=138mm,yshift=-95.00pt]current page.north west) {\textcolor{SNink}{\textbf{错项排除 · 原子锁不自动保证公平}}\par\vspace{1.5mm}TestAndSet保证检查旧值和置位不可被插入，因而可以建立互斥；它没有规定等待者先后次序。某进程持续抢不到锁仍可能饥饿，所以``硬件原子指令保证有限等待''缺少排队、公平和持续运行等条件。};}
\expandafter\gdef\csname SNside10\endcsname{\node[anchor=north west,inner sep=0pt,text width=48mm,font=\SideFont,align=left] at ([xshift=138mm,yshift=-95.00pt]current page.north west) {\textcolor{SNink}{\textbf{课程连接 · 迁移与Cache局部性相互影响}}\par\vspace{1.5mm}把线程移到空闲核可减少排队，却可能丢掉此前核的Cache与TLB热状态。组成原理的命中率和操作系统的负载均衡在这里共同影响总时间；每核任务数量相等也不保证剩余工作量相等。};}
\expandafter\gdef\csname SNside11\endcsname{\node[anchor=north west,inner sep=0pt,text width=48mm,font=\SideFont,align=left] at ([xshift=138mm,yshift=-203.72pt]current page.north west) {\textcolor{SNink}{\textbf{不变量 · 记录型计数的符号含义}}\par\vspace{1.5mm}采用本节记录型定义时，正数表示未分配许可，负值绝对值表示等待者数量。P先减再判断，结果等于0仍已拿到最后许可；V先加再判断，结果等于0仍需唤醒。不能把整型忙等信号量的非负规则混进来。};}
\expandafter\gdef\csname SNside12\endcsname{\node[anchor=north west,inner sep=0pt,text width=48mm,font=\SideFont,align=left] at ([xshift=138mm,yshift=-95.00pt]current page.north west) {\textcolor{SNink}{\textbf{错项排除 · 满足次序还要保留并发}}\par\vspace{1.5mm}若要求C完成后E可开始，却在D完成后才通知E，虽然没有违反C先于E，却额外加上D先于E。前驱题需同时检查``不提前''和``不增加无关等待''，通知应紧随真正的前驱条件成立时发出。};}
\expandafter\gdef\csname SNside13\endcsname{\node[anchor=north west,inner sep=0pt,text width=48mm,font=\SideFont,align=left] at ([xshift=138mm,yshift=-95.00pt]current page.north west) {\textcolor{SNink}{\textbf{最小反例 · 锁住改变条件的人}}\par\vspace{1.5mm}初始空缓冲区中，消费者若先P(mutex)再P(full)，就会持锁等产品；生产者随后等同一把锁，无法放入产品。两方等待的事件互相依赖。修正时应让消费者先取得产品许可，再申请只保护单次取出的缓冲锁。\par\smallskip\begin{tikzpicture}[x=1mm,y=1mm,>=Stealth,node distance=5mm]
\node[draw=SNink!50,fill=SNblue,rounded corners=1mm,text width=34mm,align=center,inner sep=1.2mm,font=\fontsize{7.5}{9}\selectfont] (m0) at (0,0) {缓冲为空};
\node[draw=SNink!50,fill=SNblue,rounded corners=1mm,text width=34mm,align=center,inner sep=1.2mm,font=\fontsize{7.5}{9}\selectfont] (m1) at (0,-11) {消费者持锁};
\node[draw=SNink!50,fill=SNblue,rounded corners=1mm,text width=34mm,align=center,inner sep=1.2mm,font=\fontsize{7.5}{9}\selectfont] (m2) at (0,-22) {生产者等锁};
\draw[->,SNink] (m0) -- (m1);
\draw[->,SNink] (m1) -- (m2);
\draw[->,SNink] (m2) -- (m1);
\end{tikzpicture}
\par{\fontsize{7.5}{9}\selectfont 消费者持mutex等待full，生产者也需要mutex，二者无法推进。}};}
\expandafter\gdef\csname SNside14\endcsname{\node[anchor=north west,inner sep=0pt,text width=48mm,font=\SideFont,align=left] at ([xshift=138mm,yshift=-95.00pt]current page.north west) {\textcolor{SNink}{\textbf{不变量 · 许可与物理槽位有中间态}}\par\vspace{1.5mm}容量2且初始为空，生产者刚做P(empty)尚未放入时，empty＝1、full＝0，物理空槽却仍为2。许可已预留而物理状态未改变，因此不能把empty＋full＝容量写成每条指令后都成立的不变量。};}
\expandafter\gdef\csname SNside15\endcsname{\node[anchor=north west,inner sep=0pt,text width=48mm,font=\SideFont,align=left] at ([xshift=138mm,yshift=-95.00pt]current page.north west) {\textcolor{SNink}{\textbf{错项排除 · 原子锁不自动保证公平}}\par\vspace{1.5mm}TestAndSet保证检查旧值和置位不可被插入，因而可以建立互斥；它没有规定等待者先后次序。某进程持续抢不到锁仍可能饥饿，所以``硬件原子指令保证有限等待''缺少排队、公平和持续运行等条件。};}
\expandafter\gdef\csname SNside16\endcsname{\node[anchor=north west,inner sep=0pt,text width=48mm,font=\SideFont,align=left] at ([xshift=138mm,yshift=-95.00pt]current page.north west) {\textcolor{SNink}{\textbf{不变量 · 记录型计数的符号含义}}\par\vspace{1.5mm}采用本节记录型定义时，正数表示未分配许可，负值绝对值表示等待者数量。P先减再判断，结果等于0仍已拿到最后许可；V先加再判断，结果等于0仍需唤醒。不能把整型忙等信号量的非负规则混进来。};}
\expandafter\gdef\csname SNside17\endcsname{\node[anchor=north west,inner sep=0pt,text width=48mm,font=\SideFont,align=left] at ([xshift=138mm,yshift=-95.00pt]current page.north west) {\textcolor{SNink}{\textbf{错项排除 · 满足次序还要保留并发}}\par\vspace{1.5mm}若要求C完成后E可开始，却在D完成后才通知E，虽然没有违反C先于E，却额外加上D先于E。前驱题需同时检查``不提前''和``不增加无关等待''，通知应紧随真正的前驱条件成立时发出。};}
\expandafter\gdef\csname SNside18\endcsname{\node[anchor=north west,inner sep=0pt,text width=48mm,font=\SideFont,align=left] at ([xshift=138mm,yshift=-95.00pt]current page.north west) {\textcolor{SNink}{\textbf{最小反例 · 锁住改变条件的人}}\par\vspace{1.5mm}初始空缓冲区中，消费者若先P(mutex)再P(full)，就会持锁等产品；生产者随后等同一把锁，无法放入产品。两方等待的事件互相依赖。修正时应让消费者先取得产品许可，再申请只保护单次取出的缓冲锁。\par\smallskip\begin{tikzpicture}[x=1mm,y=1mm,>=Stealth,node distance=5mm]
\node[draw=SNink!50,fill=SNblue,rounded corners=1mm,text width=34mm,align=center,inner sep=1.2mm,font=\fontsize{7.5}{9}\selectfont] (m0) at (0,0) {缓冲为空};
\node[draw=SNink!50,fill=SNblue,rounded corners=1mm,text width=34mm,align=center,inner sep=1.2mm,font=\fontsize{7.5}{9}\selectfont] (m1) at (0,-11) {消费者持锁};
\node[draw=SNink!50,fill=SNblue,rounded corners=1mm,text width=34mm,align=center,inner sep=1.2mm,font=\fontsize{7.5}{9}\selectfont] (m2) at (0,-22) {生产者等锁};
\draw[->,SNink] (m0) -- (m1);
\draw[->,SNink] (m1) -- (m2);
\draw[->,SNink] (m2) -- (m1);
\end{tikzpicture}
\par{\fontsize{7.5}{9}\selectfont 消费者持mutex等待full，生产者也需要mutex，二者无法推进。}};}
\expandafter\gdef\csname SNside19\endcsname{\node[anchor=north west,inner sep=0pt,text width=48mm,font=\SideFont,align=left] at ([xshift=138mm,yshift=-95.00pt]current page.north west) {\textcolor{SNink}{\textbf{不变量 · 许可与物理槽位有中间态}}\par\vspace{1.5mm}容量2且初始为空，生产者刚做P(empty)尚未放入时，empty＝1、full＝0，物理空槽却仍为2。许可已预留而物理状态未改变，因此不能把empty＋full＝容量写成每条指令后都成立的不变量。};}
\expandafter\gdef\csname SNside20\endcsname{\node[anchor=north west,inner sep=0pt,text width=48mm,font=\SideFont,align=left] at ([xshift=138mm,yshift=-95.00pt]current page.north west) {\textcolor{SNink}{\textbf{不变量 · 登记与首末判断一起保护}}\par\vspace{1.5mm}第一个读者申请文件权，最后一个读者释放。若只保护计数加减而不保护是否为首末读者的判断，新读者可能在旧末读者释放后直接进入，与写者重叠。保护对象是整个判定与登记过程，不能只盯住一条加法。};}
\expandafter\gdef\csname SNside21\endcsname{\node[anchor=north west,inner sep=0pt,text width=48mm,font=\SideFont,align=left] at ([xshift=138mm,yshift=-95.00pt]current page.north west) {\textcolor{SNink}{\textbf{错项排除 · 四人申请不等于四人吃饭}}\par\vspace{1.5mm}五根筷子下，每个吃饭者要占两根，最多只能有两人同时吃。room＝4限制的是开始申请到全部归还这一段内的人数，用来避免五人各拿一根；它既不是同时吃饭人数，也不直接保证所有等待者公平。};}
\expandafter\gdef\csname SNside22\endcsname{\node[anchor=north west,inner sep=0pt,text width=48mm,font=\SideFont,align=left] at ([xshift=138mm,yshift=-95.00pt]current page.north west) {\textcolor{SNink}{\textbf{最小反例 · 通知不会留下许可}}\par\vspace{1.5mm}条件变量没有等待者时发signal，随后才调用wait的线程不会获得一张已保存的通知券。相比之下，信号量V会增加可供以后P取得的许可。用``通知在前、等待在后''这一最小时间线就能排除二者作用相同的说法。};}
\expandafter\gdef\csname SNside23\endcsname{\node[anchor=north west,inner sep=0pt,text width=48mm,font=\SideFont,align=left] at ([xshift=138mm,yshift=-95.00pt]current page.north west) {\textcolor{SNink}{\textbf{不变量 · 登记与首末判断一起保护}}\par\vspace{1.5mm}第一个读者申请文件权，最后一个读者释放。若只保护计数加减而不保护是否为首末读者的判断，新读者可能在旧末读者释放后直接进入，与写者重叠。保护对象是整个判定与登记过程，不能只盯住一条加法。};}
\expandafter\gdef\csname SNside24\endcsname{\node[anchor=north west,inner sep=0pt,text width=48mm,font=\SideFont,align=left] at ([xshift=138mm,yshift=-95.00pt]current page.north west) {\textcolor{SNink}{\textbf{错项排除 · 四人申请不等于四人吃饭}}\par\vspace{1.5mm}五根筷子下，每个吃饭者要占两根，最多只能有两人同时吃。room＝4限制的是开始申请到全部归还这一段内的人数，用来避免五人各拿一根；它既不是同时吃饭人数，也不直接保证所有等待者公平。};}
\expandafter\gdef\csname SNside25\endcsname{\node[anchor=north west,inner sep=0pt,text width=48mm,font=\SideFont,align=left] at ([xshift=138mm,yshift=-95.00pt]current page.north west) {\textcolor{SNink}{\textbf{错项排除 · 有环还要看资源实例数}}\par\vspace{1.5mm}P持A等B、Q持B等A构成环；若A另有空闲实例，Q仍可拿到它并完成，归还B后P也能前进。单实例图中的环足以判断死锁，多实例时还须化简；不能只看箭头形状就跳过可用实例。};}
\expandafter\gdef\csname SNside26\endcsname{\node[anchor=north west,inner sep=0pt,text width=48mm,font=\SideFont,align=left] at ([xshift=138mm,yshift=-95.00pt]current page.north west) {\textcolor{SNink}{\textbf{最小反例 · 通知不会留下许可}}\par\vspace{1.5mm}条件变量没有等待者时发signal，随后才调用wait的线程不会获得一张已保存的通知券。相比之下，信号量V会增加可供以后P取得的许可。用``通知在前、等待在后''这一最小时间线就能排除二者作用相同的说法。};}
\expandafter\gdef\csname SNside27\endcsname{\node[anchor=north west,inner sep=0pt,text width=48mm,font=\SideFont,align=left] at ([xshift=138mm,yshift=-95.00pt]current page.north west) {\textcolor{SNink}{\textbf{不变量 · Work增加的是原持有量}}\par\vspace{1.5mm}安全检查假想先借出Need、再收回Max，净效果是Work增加Allocation。若直接增加Max，就把尚未借出的资源也算成新增资源，破坏总量守恒。算法中的每次假定完成，都应能还原成这次借出与归还。};}
\expandafter\gdef\csname SNside28\endcsname{\node[anchor=north west,inner sep=0pt,text width=48mm,font=\SideFont,align=left] at ([xshift=138mm,yshift=-95.00pt]current page.north west) {\textcolor{SNink}{\textbf{最小反例 · 拿得出仍可能不能批}}\par\vspace{1.5mm}总量4，P最多4、Q最多3，已各持1，可用2；当前可先完成Q。若再给P一个，可用只剩1，两者都还需2，安全序列消失。这个状态通过了现存量检查，却未通过试分配后的安全检查。};}
\expandafter\gdef\csname SNside29\endcsname{\node[anchor=north west,inner sep=0pt,text width=48mm,font=\SideFont,align=left] at ([xshift=138mm,yshift=-228.38pt]current page.north west) {\textcolor{SNink}{\textbf{错项排除 · Need不等于当前请求}}\par\vspace{1.5mm}银行家用未来可能达到的Need验证完成保证；死锁检测用已经提出、正在等待的Request检查当前能否推进。不安全说明最坏未来缺少保证，不等于当前已经死锁，更不能直接宣布所有未标记者都已死锁。};}
\expandafter\gdef\csname SNside30\endcsname{\node[anchor=north west,inner sep=0pt,text width=48mm,font=\SideFont,align=left] at ([xshift=138mm,yshift=-95.00pt]current page.north west) {\textcolor{SNink}{\textbf{错项排除 · 有环还要看资源实例数}}\par\vspace{1.5mm}P持A等B、Q持B等A构成环；若A另有空闲实例，Q仍可拿到它并完成，归还B后P也能前进。单实例图中的环足以判断死锁，多实例时还须化简；不能只看箭头形状就跳过可用实例。};}
\expandafter\gdef\csname SNside31\endcsname{\node[anchor=north west,inner sep=0pt,text width=48mm,font=\SideFont,align=left] at ([xshift=138mm,yshift=-95.00pt]current page.north west) {\textcolor{SNink}{\textbf{不变量 · Work增加的是原持有量}}\par\vspace{1.5mm}安全检查假想先借出Need、再收回Max，净效果是Work增加Allocation。若直接增加Max，就把尚未借出的资源也算成新增资源，破坏总量守恒。算法中的每次假定完成，都应能还原成这次借出与归还。};}
\expandafter\gdef\csname SNside32\endcsname{\node[anchor=north west,inner sep=0pt,text width=48mm,font=\SideFont,align=left] at ([xshift=138mm,yshift=-95.00pt]current page.north west) {\textcolor{SNink}{\textbf{最小反例 · 拿得出仍可能不能批}}\par\vspace{1.5mm}总量4，P最多4、Q最多3，已各持1，可用2；当前可先完成Q。若再给P一个，可用只剩1，两者都还需2，安全序列消失。这个状态通过了现存量检查，却未通过试分配后的安全检查。};}
\expandafter\gdef\csname SNside33\endcsname{\node[anchor=north west,inner sep=0pt,text width=48mm,font=\SideFont,align=left] at ([xshift=138mm,yshift=-255.46pt]current page.north west) {\textcolor{SNink}{\textbf{错项排除 · Need不等于当前请求}}\par\vspace{1.5mm}银行家用未来可能达到的Need验证完成保证；死锁检测用已经提出、正在等待的Request检查当前能否推进。不安全说明最坏未来缺少保证，不等于当前已经死锁，更不能直接宣布所有未标记者都已死锁。};}
\expandafter\gdef\csname SNside34\endcsname{\node[anchor=north west,inner sep=0pt,text width=48mm,font=\SideFont,align=left] at ([xshift=138mm,yshift=-95.00pt]current page.north west) {\textcolor{SNink}{\textbf{错项排除 · Need不等于当前请求}}\par\vspace{1.5mm}银行家用未来可能达到的Need验证完成保证；死锁检测用已经提出、正在等待的Request检查当前能否推进。不安全说明最坏未来缺少保证，不等于当前已经死锁，更不能直接宣布所有未标记者都已死锁。};}
\expandafter\gdef\csname SNside35\endcsname{\node[anchor=north west,inner sep=0pt,text width=48mm,font=\SideFont,align=left] at ([xshift=138mm,yshift=-281.83pt]current page.north west) {\textcolor{SNink}{\textbf{不变量 · Work增加的是原持有量}}\par\vspace{1.5mm}安全检查假想先借出Need、再收回Max，净效果是Work增加Allocation。若直接增加Max，就把尚未借出的资源也算成新增资源，破坏总量守恒。算法中的每次假定完成，都应能还原成这次借出与归还。};}
\expandafter\gdef\csname SNside36\endcsname{\node[anchor=north west,inner sep=0pt,text width=48mm,font=\SideFont,align=left] at ([xshift=138mm,yshift=-95.00pt]current page.north west) {\textcolor{SNink}{\textbf{最小反例 · 拿得出仍可能不能批}}\par\vspace{1.5mm}总量4，P最多4、Q最多3，已各持1，可用2；当前可先完成Q。若再给P一个，可用只剩1，两者都还需2，安全序列消失。这个状态通过了现存量检查，却未通过试分配后的安全检查。};}
\AddToShipoutPictureFG{\begin{tikzpicture}[remember picture,overlay]\ifcsname SNside\thepage\endcsname\csname SNside\thepage\endcsname\fi\end{tikzpicture}}